\documentclass[12pt,a4paper]{article}

% --- Metadata ---
\def\courseshortheader{HỌC MÁY NÂNG CAO --- MẠNG NƠRON}
\def\coverbookline{TRÍ TUỆ NHÂN TẠO}
\def\covermaintitle{Mạng nơron}
\def\coversubtitle{Lý thuyết --- Kiến trúc, hàm kích hoạt, lan truyền ngược, đa lớp}

% Shared preamble for nhom_5 (requires \courseshortheader before \input)
\usepackage[utf8]{inputenc}
\usepackage[vietnamese]{babel}
\usepackage{amsmath,amssymb,amsthm}
\usepackage{geometry}
\usepackage{enumitem}
\usepackage[most]{tcolorbox}
\usepackage{hyperref}
\usepackage{fancyhdr}
\usepackage{graphicx}
\usepackage{booktabs}
\usepackage{tikz}
\usetikzlibrary{calc,arrows.meta,decorations.pathreplacing,positioning}
\usepackage{eso-pic}
\usepackage{xcolor}

\geometry{a4paper, margin=2cm, bottom=2.5cm}
\hypersetup{colorlinks=true, linkcolor=blue!70!black, urlcolor=blue}
\setlist[itemize]{topsep=4pt, itemsep=2pt}
\setlist[enumerate]{topsep=4pt, itemsep=2pt}

\AddToShipoutPictureFG{%
  \AtPageCenter{%
    \begin{tikzpicture}[remember picture, overlay]
      \node[rotate=45, scale=1, opacity=0.06]{%
        \fontsize{60}{72}\selectfont\bfseries\sffamily Đặng Tấn Quí%
      };
    \end{tikzpicture}%
  }%
}

\pagestyle{fancy}
\fancyhf{}
\fancyhead[L]{\small\textit{\courseshortheader}}
\fancyhead[R]{\small\textit{Đặng Tấn Quí}}
\fancyfoot[C]{\thepage}
\fancyfoot[L]{\footnotesize\textcopyright\ Đặng Tấn Quí}
\fancyfoot[R]{\footnotesize SĐT: 0377\,122\,210}
\renewcommand{\headrulewidth}{0.4pt}
\renewcommand{\footrulewidth}{0.4pt}

\fancypagestyle{plain}{%
  \fancyhf{}
  \fancyfoot[C]{\thepage}
  \fancyfoot[L]{\footnotesize\textcopyright\ Đặng Tấn Quí}
  \fancyfoot[R]{\footnotesize SĐT: 0377\,122\,210}
  \renewcommand{\headrulewidth}{0pt}
  \renewcommand{\footrulewidth}{0.4pt}
}

\newtcolorbox{dongco}[1][]{%
  enhanced, breakable,
  colback=teal!4!white, colframe=teal!60!black,
  fonttitle=\bfseries, title={Động cơ học -- Tại sao cần khái niệm này?}, #1%
}
\newtcolorbox{ynghia}[1][]{%
  enhanced, breakable,
  colback=purple!3!white, colframe=purple!50!black,
  fonttitle=\bfseries, title={Ý nghĩa \& Trực giác}, #1%
}
\newtcolorbox{ungdung}[1][]{%
  enhanced, breakable,
  colback=olive!5!white, colframe=olive!60!black,
  fonttitle=\bfseries, title={Ứng dụng thực tế}, #1%
}
\newtcolorbox{dinhnghia}[1][]{%
  enhanced, breakable,
  colback=blue!3!white, colframe=blue!60!black,
  fonttitle=\bfseries, title={Định nghĩa}, #1%
}
\newtcolorbox{dinhly}[1][]{%
  enhanced, breakable,
  colback=green!3!white, colframe=green!50!black,
  fonttitle=\bfseries, title={Định lý}, #1%
}
\newtcolorbox{tinhchat}[1][]{%
  enhanced, breakable,
  colback=yellow!5!white, colframe=orange!50!black,
  fonttitle=\bfseries, title={Tính chất}, #1%
}
\newtcolorbox{phuongphap}[1][]{%
  enhanced, breakable,
  colback=gray!5!white, colframe=gray!50!black,
  fonttitle=\bfseries, title={Phương pháp}, #1%
}
\newtcolorbox{luuy}[1][]{%
  enhanced, breakable,
  colback=red!3!white, colframe=red!50!black,
  fonttitle=\bfseries, title={Lưu ý}, #1%
}
\newtcolorbox{congthuc}[1][]{%
  enhanced, breakable,
  colback=cyan!3!white, colframe=cyan!50!black,
  fonttitle=\bfseries, title={Công thức}, #1%
}
\newtcolorbox{vidu}[1][]{%
  enhanced, breakable,
  colback=violet!3!white, colframe=violet!50!black,
  fonttitle=\bfseries, title={Ví dụ}, #1%
}
\newtcolorbox{phanvidu}[1][]{%
  enhanced, breakable,
  colback=red!2!white, colframe=red!40!black,
  fonttitle=\bfseries, title={Phản ví dụ}, #1%
}
\newtcolorbox{tongket}[1][]{%
  enhanced, breakable,
  colback=blue!4!white, colframe=blue!70!black,
  fonttitle=\bfseries, title={Tổng kết chương}, #1%
}


\usepackage{tabularx}
\usepackage{array}
\usepackage{float}
\usepackage{amsmath}

\graphicspath{{images/neural-networks/}}

\newcommand{\nnimg}[2][0.88]{%%
  \def\nnfile{images/neural-networks/#2}%%
  \IfFileExists{\nnfile}{%%
    \begin{center}\includegraphics[width=#1\linewidth]{#2}\end{center}%%
  }{%%
    \begin{luuy}[title={Thiếu hình minh họa}]%%
      Không tìm thấy \path{\nnfile}. Nếu cần, tải PNG tương ứng từ MLCC.%%
    \end{luuy}%%
  }%%
}

\begin{document}

\thispagestyle{empty}
\begin{tikzpicture}[remember picture, overlay]
  \fill[blue!80!black] (current page.south west) rectangle (current page.north east);
  \fill[blue!60!cyan, opacity=0.8]
    (current page.south west) -- (current page.north west) --
    ([xshift=-3cm]current page.north east) -- (current page.south east) -- cycle;
  \foreach \r/\o in {6/0.05, 4.5/0.08, 3/0.12} {
    \fill[white, opacity=\o] ([xshift=2cm, yshift=-2cm]current page.north east) circle (\r cm);
  }
  \draw[white, line width=2pt] ([yshift=-5cm]current page.north west) ++(2cm,0) -- ++(17cm,0);
  \draw[white, line width=2pt] ([yshift=-16cm]current page.north west) ++(2cm,0) -- ++(17cm,0);
  \node[anchor=west, white, font=\fontsize{28}{34}\bfseries\selectfont]
    at ([xshift=3cm, yshift=-7.5cm]current page.north west) {\coverbookline};
  \node[anchor=west, white, font=\fontsize{20}{26}\selectfont]
    at ([xshift=3cm, yshift=-9cm]current page.north west) {\covermaintitle};
  \node[anchor=west, white, font=\fontsize{16}{22}\selectfont]
    at ([xshift=3cm, yshift=-10.2cm]current page.north west) {\coversubtitle};
  \node[anchor=west, white, font=\Large\bfseries]
    at ([xshift=3cm, yshift=-13cm]current page.north west) {Biên soạn:};
  \node[anchor=west, yellow!80!white, font=\fontsize{24}{30}\bfseries\selectfont]
    at ([xshift=3cm, yshift=-14.5cm]current page.north west) {ĐẶNG TẤN QUÍ};
  \node[anchor=west, white!90, font=\large]
    at ([xshift=3cm, yshift=-18cm]current page.north west) {SĐT: 0377\,122\,210};
  \node[anchor=south, white!80, font=\small]
    at ([yshift=1.5cm]current page.south) {Tài liệu có bản quyền --- Cấm sao chép khi chưa được sự đồng ý của tác giả};
  \node[anchor=south east, white, font=\Large\bfseries]
    at ([xshift=-2cm, yshift=3cm]current page.south east) {\the\year};
\end{tikzpicture}

\newpage
\tableofcontents
\newpage

%==============================================================================
\section*{Lời nói đầu}
\addcontentsline{toc}{section}{Lời nói đầu}
%==============================================================================

\begin{ynghia}[title={Nguồn và cách đọc}]
Tài liệu biên soạn và dịch từ \emph{Neural networks} (Google Machine Learning Crash Course).
\end{ynghia}

\begin{luuy}[title={Điều kiện tiên quyết}]
\begin{itemize}
  \item Introduction to ML, Linear regression, Logistic regression, Classification.
  \item Numerical data, Categorical data.
  \item Datasets, Generalization, and Overfitting.
\end{itemize}
\end{luuy}

\begin{phuongphap}[title={Mục tiêu học tập}]
\begin{itemize}
  \item Hiểu động cơ dùng mạng nơron cho bài toán phi tuyến.
  \item Nắm khái niệm nút, lớp ẩn, hàm kích hoạt.
  \item Nắm trực giác suy luận (inference) trong mạng nơron nhiều lớp.
  \item Nắm ý tưởng lan truyền ngược, vấn đề gradient mất/ nổ, dropout.
  \item Hiểu phân loại đa lớp kiểu one-vs-all và softmax (one-vs-one).
\end{itemize}
\end{phuongphap}

\newpage

%==============================================================================
\section{Giới thiệu mạng nơron}
%==============================================================================

\begin{dongco}[title={Bài toán phi tuyến}]
Trong module dữ liệu danh mục, bạn đã thấy bài toán phân loại phi tuyến: điểm xanh/da cam nằm xen kẽ theo hình vòng tròn hoặc hình chữ thập; đường thẳng không thể chia chúng sạch sẽ.
\end{dongco}

\begin{vidu}[title={Feature cross cho bài toán đơn giản}]
Với bài toán ``chữ thập'' đơn giản, thêm feature cross $x_1x_2$ cho phép mô hình tuyến tính học đường cong siêu bậc tách hai lớp.
\end{vidu}

\begin{vidu}[title={Bài toán vòng tròn khó hơn}]
Với tập dữ liệu ``vòng tròn trong vòng tròn'', việc đoán trước feature cross phù hợp trở nên khó và tốn công thử nghiệm.
\end{vidu}

\begin{dinhnghia}[title={Mạng nơron}]
\href{https://developers.google.com/machine-learning/glossary\#neural_network}{Mạng nơron} là họ kiến trúc mô hình được thiết kế để học \href{https://developers.google.com/machine-learning/glossary\#nonlinear}{mối quan hệ phi tuyến} trong dữ liệu, bằng cách xếp chồng nhiều tầng phép biến đổi tuyến tính + phi tuyến.
\end{dinhnghia}

\begin{tongket}
Các phần sau lần lượt trình bày: (1) nút và lớp ẩn, (2) hàm kích hoạt, (3) lan truyền ngược và regularization, (4) bài tập tương tác, (5) phân loại đa lớp.
\end{tongket}

%==============================================================================
\section{Nút và lớp ẩn}
%==============================================================================

\begin{dinhnghia}[title={Mô hình tuyến tính dạng mạng}]
Mô hình tuyến tính quen thuộc $y' = b + w_1x_1 + w_2x_2 + w_3x_3$ có thể vẽ thành ``mạng'' với 3 nút đầu vào (xanh) và 1 nút đầu ra (xanh lá): mỗi cạnh mang một trọng số $w_i$, cộng thêm bias $b$.
\end{dinhnghia}

\begin{phuongphap}[title={Bài tập 1 (tương tác trong MLCC)}]
Widget cho phép thay đổi các $x_i$, $w_i$, $b$ và xem lại cách tính $y'$ (hàm Linear chỉ trả về đầu vào). Việc chỉnh tham số không làm thay đổi việc mô hình vẫn là \textbf{tuyến tính} theo $x_i$.
\end{phuongphap}

\begin{dinhnghia}[title={Lớp ẩn và nơron}]
\begin{itemize}
  \item \href{https://developers.google.com/machine-learning/glossary\#input-layer}{Lớp đầu vào}: chứa các đặc trưng $x_i$.
  \item \href{https://developers.google.com/machine-learning/glossary\#hidden_layer}{Lớp ẩn}: lớp nằm giữa đầu vào và đầu ra; mỗi nút gọi là \href{https://developers.google.com/machine-learning/glossary\#neuron}{nơron}.
  \item Giá trị từng nơron ẩn: tổng có trọng số của đầu vào từ lớp trước + bias riêng, giống hệt mô hình tuyến tính.
\end{itemize}
\end{dinhnghia}

\begin{ynghia}[title={Thêm lớp ẩn có giúp phi tuyến?}]
Nếu mọi phép biến đổi giữa các lớp đều \textbf{tuyến tính} (nhân + cộng), thì xếp chồng nhiều lớp vẫn chỉ là \textbf{một phép tuyến tính tương đương}. Do đó, chỉ thêm lớp ẩn mà không thêm phi tuyến thì mô hình vẫn không học được phi tuyến.
\end{ynghia}

\begin{luuy}[title={Kiểm tra hiểu biết --- Số tham số}]
Ví dụ mạng 3 đầu vào, 1 lớp ẩn 4 nơron, 1 đầu ra:
\begin{itemize}
  \item 4 nơron ẩn: mỗi nơron có 3 trọng số + 1 bias $\Rightarrow 16$ tham số.
  \item 1 nơron đầu ra: 4 trọng số (mỗi nơron ẩn một trọng số) + 1 bias $\Rightarrow 5$ tham số.
  \item Tổng cộng: 21 tham số.
\end{itemize}
Mặc dù nhiều lớp và tham số hơn, mọi phép tính vẫn tuyến tính $\Rightarrow$ chưa học được phi tuyến.
\end{luuy}

%==============================================================================
\section{Hàm kích hoạt}
%==============================================================================

\begin{dongco}[title={Làm sao đưa phi tuyến vào mạng?}]
Ghép tuyến tính sau tuyến tính vẫn tuyến tính. Cần chèn \textbf{phép biến đổi phi tuyến} giữa các lớp để mạng học quan hệ phức tạp.
\end{dongco}

\begin{vidu}[title={So sánh với logistic regression}]
Trong hồi quy logistic, ta đưa đầu ra tuyến tính $z$ qua hàm sigmoid để được xác suất $y'$. Tương tự, trong mạng nơron, ta đưa tổng có trọng số qua \href{https://developers.google.com/machine-learning/glossary\#activation_function}{\textbf{hàm kích hoạt}} trước khi chuyển sang lớp kế tiếp.
\end{vidu}

\begin{dinhnghia}[title={Hàm kích hoạt}]
Hàm kích hoạt $\sigma$ là hàm phi tuyến áp dụng lên đầu ra tuyến tính của mỗi nơron:
\[
  \text{nơron}(\boldsymbol x) = \sigma(\boldsymbol w \cdot \boldsymbol x + b)
\]
Xếp chồng nhiều lớp với hàm kích hoạt $\Rightarrow$ mô hình hóa được quan hệ phi tuyến phức tạp.
\end{dinhnghia}

\subsection{Ba hàm kích hoạt phổ biến}

\begin{congthuc}[title={Sigmoid}]
\[
  \sigma(x) = \frac{1}{1 + e^{-x}}
\]
Đầu ra trong $(0,1)$. Hay dùng ở đầu ra cho xác suất.
\end{congthuc}

\begin{congthuc}[title={tanh}]
\[
  \sigma(x) = \tanh(x)
\]
Đầu ra trong $(-1,1)$. Hình chữ S giống sigmoid nhưng dốc hơn.
\end{congthuc}

\begin{congthuc}[title={ReLU}]
\[
  \sigma(x) = \max(0,x)
\]
Đưa mọi giá trị âm về 0, giữ nguyên giá trị dương.
\end{congthuc}

\begin{tinhchat}[title={So sánh sigmoid, tanh, ReLU}]
\begin{itemize}
  \item Sigmoid, tanh trơn, nhưng dễ gặp \href{https://developers.google.com/machine-learning/glossary\#vanishing-gradient-problem}{vanishing gradient} khi $|x|$ lớn.
  \item ReLU đơn giản, rẻ tính toán, thường giúp giảm vấn đề gradient mất.
\end{itemize}
\end{tinhchat}

\begin{luuy}[title={Tóm tắt kiến trúc mạng nơron}]
Một mạng nơron điển hình gồm:
\begin{itemize}
  \item Nhiều lớp nơron (input, ẩn, output).
  \item Trọng số + bias giữa các lớp.
  \item Hàm kích hoạt áp dụng trên mỗi nơron (có thể khác nhau giữa các lớp, nhưng thường cùng loại trong một lớp).
\end{itemize}
\end{luuy}

%==============================================================================
\section{Huấn luyện bằng lan truyền ngược}
%==============================================================================

\begin{dinhnghia}[title={Backpropagation}]
\href{https://developers.google.com/machine-learning/glossary\#backpropagation}{Lan truyền ngược} là thuật toán chuẩn để tính gradient trong mạng nơron nhiều lớp, giúp áp dụng gradient descent cho toàn mạng.
\end{dinhnghia}

\begin{ungdung}[title={Video trực quan}]
\url{https://www.youtube.com/watch?v=ovCyBmlrZGg}
\end{ungdung}

\subsection{Một số hiện tượng khi huấn luyện}

\begin{tinhchat}[title={Vanishing gradient}]
\begin{itemize}
  \item Ở các lớp dưới (gần đầu vào), gradient có thể rất nhỏ do là tích của nhiều số nhỏ.
  \item Khi gradient gần 0, lớp đó học rất chậm hoặc không học được.
  \item Hàm kích hoạt như ReLU giúp giảm bớt hiện tượng này.
\end{itemize}
\end{tinhchat}

\begin{tinhchat}[title={Exploding gradient}]
\begin{itemize}
  \item Nếu trọng số lớn, gradient ở lớp dưới là tích của nhiều số lớn $\Rightarrow$ \textbf{gradient nổ}.
  \item Có thể làm việc học không hội tụ.
  \item Khắc phục: batch normalization, giảm learning rate.
\end{itemize}
\end{tinhchat}

\begin{tinhchat}[title={Dead ReLU units}]
\begin{itemize}
  \item Khi tổng có trọng số vào ReLU $<0$, đầu ra = 0 và gradient qua nút đó = 0.
  \item Nút có thể ``chết'' (luôn ra 0) nếu không còn gradient để cập nhật.
  \item Khắc phục: giảm learning rate; dùng biến thể ReLU (LeakyReLU, v.v.).
\end{itemize}
\end{tinhchat}

\subsection{Dropout regularization}

\begin{dinhnghia}[title={Dropout}]
\href{https://developers.google.com/machine-learning/glossary\#dropout_regularization}{Dropout} là kỹ thuật regularization trong đó một tỉ lệ nút được ``tắt'' (output=0) ngẫu nhiên trong một bước cập nhật.
\begin{itemize}
  \item Tỉ lệ 0,0: không dropout.
  \item Tỉ lệ 1,0: tắt hết, không học được gì.
  \item Thực tế: giá trị trung gian (ví dụ 0,2--0,5).
\end{itemize}
\end{dinhnghia}

%==============================================================================
\section{Bài tập tương tác}
%==============================================================================

\begin{phuongphap}[title={Khám phá mạng nơron bằng widget}]
MLCC cung cấp widget cho phép:
\begin{itemize}
  \item Thay đổi số lớp ẩn, số nơron, learning rate, hàm kích hoạt.
  \item Quan sát giá trị từng nơron và loss train/test khi huấn luyện.
\end{itemize}
\end{phuongphap}

\begin{ynghia}[title={Các bài tập chính}]
\begin{itemize}
  \item Ảnh hưởng của giá trị đầu vào và trọng số tới output khi dùng ReLU.
  \item Ảnh hưởng của thêm lớp ẩn, thêm nơron vào mạng.
  \item Bài toán ``vòng tròn trong vòng tròn'': cấu hình mạng để loss train/test < 0,2.
\end{itemize}
\end{ynghia}

\begin{luuy}[title={Một cấu hình mẫu cho bài toán phi tuyến}]
Một cấu hình đạt loss < 0,2 trên train/test:
\begin{itemize}
  \item 1 lớp ẩn với 3 nơron.
  \item Learning rate 0,01.
  \item Hàm kích hoạt ReLU.
\end{itemize}
\end{luuy}

%==============================================================================
\section{Phân loại đa lớp với mạng nơron}
%==============================================================================

\begin{dinhnghia}[title={Bài toán đa lớp}]
\href{https://developers.google.com/machine-learning/glossary\#multi-class}{Phân loại đa lớp}: mỗi ví dụ thuộc một trong nhiều lớp (ví dụ loại hoa, loại máy bay, loại quả).
\end{dinhnghia}

\subsection{One-vs-all}

\begin{phuongphap}[title={Nhiều bộ phân loại nhị phân}]
\begin{itemize}
  \item Với $N$ lớp, huấn luyện $N$ bộ phân loại nhị phân ``lớp i vs không phải lớp i''.
  \item Ví dụ: ảnh quả: ``apple vs not-apple'', ``orange vs not-orange'', v.v.
\end{itemize}
\end{phuongphap}

\begin{vidu}[title={One-vs-all với mạng nơron}]
Mạng nơron nhiều lớp có thể thực hiện one-vs-all bằng cách có \textbf{một nút đầu ra cho mỗi lớp}, áp dụng sigmoid độc lập trên từng nút. Mỗi nút cho xác suất ``là lớp i'' so với ``không phải i''.
\end{vidu}

\subsection{Softmax (one-vs-one)}

\begin{dongco}[title={Vấn đề tổng xác suất}]
Trong one-vs-all, các xác suất đầu ra không nhất thiết tổng bằng 1. Khi cần xác suất tương đối giữa các lớp (ví dụ ``apple vs orange vs pear vs grape''), dùng \textbf{softmax}.
\end{dongco}

\begin{dinhnghia}[title={Softmax}]
Với véc-tơ điểm $z_j = \boldsymbol w_j^T\boldsymbol x + b_j$ cho mỗi lớp $j$ trong tập lớp $K$:
\[
  p(y=j\mid\boldsymbol x) = \frac{e^{z_j}}{\sum_{k\in K} e^{z_k}}
\]
Mọi xác suất $p(y=j\mid\boldsymbol x)$ đều trong $(0,1)$ và \textbf{tổng bằng 1}.
\end{dinhnghia}

\begin{tinhchat}[title={Biến thể softmax}]
\begin{itemize}
  \item \textbf{Full softmax}: tính xác suất cho mọi lớp --- phù hợp khi số lớp nhỏ/vừa.
  \item \textbf{Candidate sampling}: tính xác suất cho mọi nhãn dương và một mẫu con nhãn âm --- hiệu quả hơn khi số lớp rất lớn.
\end{itemize}
\end{tinhchat}

\begin{luuy}[title={Giới hạn của softmax}]
Softmax giả định \textbf{mỗi ví dụ thuộc đúng một lớp}. Nếu ví dụ có thể thuộc nhiều lớp cùng lúc (đa nhãn), cần dùng \textbf{nhiều hồi quy logistic nhị phân} thay vì softmax.
\end{luuy}

\begin{tongket}[title={Tổng kết chương --- Mạng nơron}]
\begin{itemize}
  \item Mạng nơron = nhiều lớp tuyến tính + hàm kích hoạt phi tuyến.
  \item Lan truyền ngược cho phép lan truyền gradient qua nhiều lớp; cần chú ý gradient mất/nổ, dead ReLU, dropout.
  \item Softmax mở rộng logistic cho phân loại đa lớp, softmax layer cho output tổng xác suất = 1.
\end{itemize}
\end{tongket}

\end{document}

